% ============================================================================
% Fichier  : partie3/P03_C10_cadre_juridique_cameroun_afrique.tex
% Titre    : Cadre Juridique et Réglementaire Camerounais et Africain
% Auteur   : MINKA MI NGUIDJOÏ Thierry Emmanuel
% Version  : 1.0 -- octobre 2026
% ============================================================================
% Chapitre juridique. Sources et degré de vérification :
%  - Loi n° 2010/012 : texte officiel anglais consulté (sections citées
%    textuellement).
%  - Règlement COBAC R-2016/04 : articles 1 et 2 relevés sur une base
%    réglementaire en ligne ; à collationner avec la version officielle.
%  - Loi n° 2024/017 : présentations publiées (presse spécialisée) ; aucun
%    article cité textuellement. A VERIFIER au JO.
%  - AUSCGIE (OHADA) : numéros d'articles d'après la doctrine et la
%    brochure OHADA de 2014 ; aucun article cité textuellement.
%  - Convention de Malabo : entrée en vigueur vérifiée ; statut du Cameroun
%    non établi.
% ============================================================================

\chapter{Cadre Juridique et Réglementaire Camerounais et Africain}
\label{chap:droit}
\chaptermark{Cadre juridique camerounais et africain}

\epigraph{%
    « Nul n'est censé ignorer la loi, pas même le DSI. »%
}{Adage revisité}

\niveauChapitre{Fondamental}{%
    Maîtrise du Chapitre~\ref{chap:normes} (normes, certification,
    conformité).%
}

\begin{boxobjectifs}
\begin{itemize}
    \item Identifier les obligations des lois n°~2010/012 et 2010/013 qui
          touchent la gouvernance des SI.
    \item Exposer le rôle de l'ANTIC et des autres régulateurs (ART,
          COBAC).
    \item Appliquer les exigences de la loi de 2024 sur les données
          personnelles à un système d'information.
    \item Situer les responsabilités des dirigeants dans le droit OHADA.
    \item Identifier les obligations de contrôle interne et de risque
          opérationnel des banques en zone CEMAC.
    \item Situer la Convention de Malabo.
\end{itemize}
\end{boxobjectifs}

\bigskip

% ============================================================================
\section*{Question directrice : quelles obligations ont été méconnues~?}
% ============================================================================

Trois semaines après la panne, le conseil d'administration de la BANQUE
SANAGA reçoit une note de son avocat. Elle commence ainsi~: « La banque
risque d'avoir à répondre de ses manquements devant trois autorités au
moins~: la COBAC, au titre du contrôle interne~; l'ANTIC, au titre de la
sécurité de ses systèmes d'information~; et, le cas échéant, les
juridictions, au titre de la responsabilité de ses dirigeants. » À
Bafoussam, NOUN PAY reçoit au même moment une demande d'information de
l'Autorité de protection des données à caractère personnel sur son
modèle de notation.

Ce chapitre dresse la carte de ces obligations. Il suit l'ordre du
modèle~: le fait, la règle, le texte, l'analyse. Il présente d'abord le
panorama des textes (section~\ref{sec:c10-panorama}), puis la loi sur la
cybersécurité (section~\ref{sec:c10-2010-012}), les régulateurs
(section~\ref{sec:c10-regulateurs}), les données personnelles
(section~\ref{sec:c10-donnees}), la réglementation bancaire
(section~\ref{sec:c10-cobac}), le droit des sociétés
(section~\ref{sec:c10-ohada}) et le cadre africain
(section~\ref{sec:c10-afrique}). Il se termine par les erreurs de
qualification les plus fréquentes (section~\ref{sec:c10-erreurs}).

\begin{boxalert}
\textbf{Règle de lecture.} Ce chapitre cite textuellement les seuls
articles dont le texte a pu être vérifié. Pour les autres, il en décrit la
substance d'après des sources secondaires identifiées, sans guillemets.
Avant tout usage professionnel (avis juridique, rapport au régulateur),
reportez-vous au texte publié au Journal officiel.
\end{boxalert}

\separateur

% ============================================================================
\section{Panorama des Textes}
\label{sec:c10-panorama}
% ============================================================================

\begin{tableaucours}{L{3.6cm}L{3.8cm}L{2.4cm}L{2.4cm}}{%
    \tableauHeader{Texte} &
    \tableauHeader{Objet} &
    \tableauHeader{Autorité} &
    \tableauHeader{Cas du cours}
}
Loi n°~2010/012 du 21~décembre~2010 & Cybersécurité et cybercriminalité
& ANTIC & Tous \\
Loi n°~2010/013 du 21~décembre~2010, modifiée en 2015 & Communications
électroniques & ART, MINPOSTEL & NOUN PAY \\
Loi n°~2024/017 du 23~décembre~2024 & Protection des données à caractère
personnel & Autorité de protection des données & NOUN PAY, SANAGA \\
Règlement COBAC R-2016/04 & Contrôle interne des établissements de crédit
& COBAC & SANAGA \\
Acte uniforme OHADA (révisé en 2014) & Sociétés commerciales,
responsabilité des dirigeants & Juridictions & SANAGA, KRIBILOG \\
Lois n°~2017/010 et 2017/011 & Établissements et entreprises publics &
Tutelle & ANPP \\
Loi n°~2018/012 & Régime financier de l'État & MINFI, Chambre des comptes &
ANPP \\
Convention de Malabo (2014) & Cybersécurité et données personnelles en
Afrique & Union africaine & · \\
\end{tableaucours}
\vspace{-0.8em}
\begin{center}{\small\itshape Tableau~10.1. Les principaux textes
applicables à la gouvernance des SI au Cameroun}\end{center}

% ============================================================================
\section{La Loi sur la Cybersécurité}
\label{sec:c10-2010-012}
% ============================================================================

\subsection{L'ANTIC, régulateur de la sécurité}

La loi n°~2010/012 confie la régulation de la sécurité électronique à
l'Agence nationale des technologies de l'information et de la
communication (\sigle{ANTIC}), instituée par la loi sur les
communications électroniques~\cite{loi2010012}.

\begin{boxjuris}
\textbf{\loicam{n°~2010/012}, article~7, alinéa~2} (version anglaise
officielle)~:

\medskip
\textit{« The Agency \textelp{} shall be responsible for the regulation,
control and monitoring of activities related to the security of
electronic communication networks, information systems, and electronic
certification on behalf of the State. »}

\medskip
Soit~: l'Agence est chargée, pour le compte de l'État, de la régulation,
du contrôle et du suivi des activités liées à la sécurité des réseaux de
communications électroniques, des systèmes d'information et de la
certification électronique. Parmi ses missions, la loi énumère notamment
le contrôle des activités de sécurité des systèmes d'information,
la régularité et l'efficacité des audits de sécurité, et la veille sur
les risques informatiques et la cybercriminalité.
\end{boxjuris}

\subsection{Les obligations des exploitants de systèmes d'information}

La loi impose aux exploitants de systèmes d'information
(\textit{operators of information systems}) un ensemble d'obligations qui
recoupent directement plusieurs objectifs de COBIT.

\begin{boxjuris}
\textbf{\loicam{n°~2010/012}, article~26, alinéa~1}~:

\medskip
\textit{« Operators of information systems shall take every technical and
administrative measure to ensure the security of services offered. To
this end, they shall have standardized systems enabling them to at all
times identify, assess, process or manage any risk relating to the
security of the information systems of the services provided directly or
indirectly. »}

\medskip
\textbf{Article~29, alinéa~1}~: \textit{« Operators of information systems
shall be bound to conserve the connection and traffic data of their
information systems for a period of 10 (ten) years. »}

\medskip
\textbf{Article~30, alinéa~1}~: \textit{« Operators of information systems
shall assess and revise their security systems and, where necessary, make
the appropriate modifications to their security practices, measures and
techniques according to technological change. »}

\medskip
\textbf{Analyse.} L'article~26 impose plus qu'une sécurité technique~: il
exige des \emph{dispositifs normalisés} permettant d'identifier, d'évaluer
et de traiter en permanence les risques de sécurité. C'est, en termes de
gouvernance, l'exigence d'un processus de gestion des risques (APO12) et
d'un système de management de la sécurité (APO13), dont l'ISO/IEC~27001
fournit précisément le modèle. L'article~30 consacre l'amélioration
continue. L'article~29 impose une durée de conservation des données de
connexion et de trafic de dix ans, qui pèse sur l'architecture et le
budget de stockage.
\end{boxjuris}

\begin{tableaucours}{L{2.4cm}L{5.4cm}L{4.4cm}}{%
    \tableauHeader{Article} &
    \tableauHeader{Obligation} &
    \tableauHeader{Objectifs COBIT concernés}
}
13 et 32 & Audit de sécurité obligatoire et périodique & MEA03, MEA04 \\
26 & Mesures de sécurité, dispositifs normalisés de gestion des risques &
APO12, APO13, DSS05 \\
29 & Conservation des données de connexion et de trafic pendant dix ans &
APO14, DSS01, DSS05 \\
30 & Évaluation et révision des dispositifs de sécurité & APO13, MEA01 \\
31 & Disponibilité des contenus et des données & DSS04, BAI04 \\
42 & Confidentialité des informations acheminées & DSS05, APO14 \\
\end{tableaucours}
\vspace{-0.8em}
\begin{center}{\small\itshape Tableau~10.2. Obligations de la loi
n°~2010/012 et objectifs COBIT correspondants}\end{center}

\subsection{Les sanctions et la responsabilité des personnes morales}

\begin{boxjuris}
\textbf{\loicam{n°~2010/012}, article~64, alinéas~1 à~3}~:

\medskip
\textit{« (1) Corporate bodies shall be criminally liable for offences
committed on their account by their management structures. (2) The
criminal liability of corporate bodies shall not preclude that of natural
persons who commit such offences or are accomplices. (3) The penalties to
be meted out on defaulting corporate bodies shall be fines of from
5~000~000 (five million) to 50~000~000 (fifty million) CFA francs. »}

\medskip
\textbf{Analyse.} La personne morale répond pénalement des infractions
commises pour son compte par ses organes de direction, sans que cela
exonère les personnes physiques. L'article prévoit en outre des peines
complémentaires (dissolution, interdiction d'activité, fermeture
d'établissement) dans les cas les plus graves. L'article~63 punit par
ailleurs d'emprisonnement et d'amende le dirigeant de droit ou de fait de
certains prestataires (fournisseurs d'accès et d'hébergement) qui ne
conserve pas les données exigées par les articles~25 et~29. La
gouvernance des SI n'est donc pas qu'une question de performance~: elle
engage la responsabilité pénale de l'organisation et, dans certains cas,
celle de ses dirigeants.
\end{boxjuris}

\begin{boxmetier}{DSI}
\textbf{La conservation des données pendant dix ans}
\begin{enumerate}
    \item Établir si l'organisation est « exploitant de système
          d'information » au sens de la loi~: une banque, un opérateur, une
          administration le sont presque toujours.
    \item Recenser les données de connexion et de trafic produites
          (journaux d'accès, journaux réseau, traces des applications).
    \item Définir l'architecture de conservation~: durée, intégrité (les
          journaux ne doivent pas pouvoir être modifiés), confidentialité,
          coût sur dix ans.
    \item Concilier cette obligation avec le principe de limitation de la
          conservation des données personnelles (section~\ref{sec:c10-donnees})~:
          les deux textes doivent être lus ensemble, avec un avis
          juridique.
\end{enumerate}
\end{boxmetier}

% ============================================================================
\section{Les Régulateurs}
\label{sec:c10-regulateurs}
% ============================================================================

La gouvernance des SI d'une organisation camerounaise s'exerce sous le
regard de plusieurs autorités, qu'il faut savoir distinguer.

\begin{tableaucours}{L{2.6cm}L{4.6cm}L{5.0cm}}{%
    \tableauHeader{Autorité} &
    \tableauHeader{Fondement} &
    \tableauHeader{Ce qui intéresse la gouvernance des SI}
}
MINPOSTEL & Ministère chargé des télécommunications & Politique du
secteur~; destinataire des rapports d'audit de sécurité (loi
n°~2010/012, art.~32) \\
ANTIC & Lois n°~2010/012 et 2010/013 & Sécurité des systèmes
d'information, audits, certification électronique, veille \\
ART & Loi n°~2010/013 modifiée & Régulation des communications
électroniques et des opérateurs \\
Autorité de protection des données & Loi n°~2024/017 & Traitements de
données personnelles, plaintes, sanctions \\
COBAC & Réglementation bancaire de la CEMAC & Contrôle interne et risques
des établissements de crédit, dont le risque informatique \\
Chambre des comptes & Contrôle des finances publiques & Audit des
établissements et entreprises publics, y compris de leurs SI \\
\end{tableaucours}
\vspace{-0.8em}
\begin{center}{\small\itshape Tableau~10.3. Les autorités de la gouvernance
des SI au Cameroun}\end{center}

La loi n°~2010/013 régit les communications électroniques~; elle a été
modifiée et complétée par la loi n°~2015/006 du 20~avril~2015. Elle
institue l'Agence de régulation des télécommunications (\sigle{ART}) et
l'ANTIC. Ses dispositions intéressent surtout les opérateurs et les
fournisseurs de services de communications électroniques~: une
entreprise de paiement mobile comme NOUN PAY, qui s'appuie sur les
réseaux des opérateurs, doit en connaître les grandes règles. Le recueil
des textes publié par l'ART en donne la version consolidée.

% ============================================================================
\section{La Protection des Données Personnelles}
\label{sec:c10-donnees}
% ============================================================================

\subsection{De la loi de 2010 à la loi de 2024}

La loi n°~2010/012 posait déjà un principe général.

\begin{boxjuris}
\textbf{\loicam{n°~2010/012}, articles~41 et~42}~:

\medskip
\textit{« 41. Every individual shall have the right to the protection of
their privacy. \textelp{} 42. The confidentiality of information
channelled through electronic communication and information systems
networks, including traffic data, shall be ensured by operators of
electronic communication and networks information systems. »}

\medskip
\textbf{Portée.} Ces articles affirment le droit au respect de la vie
privée et l'obligation de confidentialité des exploitants. Ils ne
constituaient pas, à eux seuls, un régime complet de protection des
données personnelles~: c'est l'objet de la loi de 2024.
\end{boxjuris}

\subsection{La loi n° 2024/017}

La loi n°~2024/017 du 23~décembre~2024 relative à la protection des
données à caractère personnel a doté le Cameroun d'un régime complet.
D'après les présentations publiées~\cite{cio2026donnees}, elle a accordé
aux organisations un délai de dix-huit mois pour se mettre en conformité,
et s'applique pleinement depuis le 23~juin~2026. Elle institue une
Autorité de protection des données à caractère personnel, organisme
public indépendant, encadre strictement le transfert des données hors du
territoire national, consacre une part importante de ses dispositions aux
obligations des responsables du traitement et de leurs sous-traitants,
et prévoit des sanctions pénales.

\begin{boxalert}
\textbf{À vérifier avant usage.} Les numéros d'articles, les définitions
exactes (responsable du traitement, sous-traitant, données sensibles),
les conditions du transfert des données hors du Cameroun et l'éventuelle
obligation de désigner un délégué à la protection des données doivent
être lus dans le texte publié et ses décrets d'application. Les
présentations disponibles divergent sur certains de ces points.
\end{boxalert}

\subsection{Ce que la loi change pour la gouvernance des SI}

Quelle que soit la lettre exacte de ses articles, une loi de protection
des données de cette génération a des conséquences directes sur la
gouvernance des SI~:
\begin{itemize}
    \item \textbf{un inventaire des traitements}~: savoir quelles données
          personnelles sont traitées, pourquoi, où et par qui (APO14)~;
    \item \textbf{la sécurité des traitements}~: mesures techniques et
          organisationnelles proportionnées aux risques (APO13, DSS05)~;
    \item \textbf{la maîtrise des sous-traitants}~: un hébergeur ou un
          fournisseur de nuage traite les données pour le compte de
          l'organisation, qui en reste responsable (APO10)~;
    \item \textbf{la localisation des données}~: tout transfert hors du
          territoire doit être examiné au regard de la loi (BAI03, APO10)~;
    \item \textbf{la conformité démontrable}~: registres, analyses
          d'impact, preuves (MEA03).
\end{itemize}

\begin{boxscenario}{Opération NOUN. La demande de l'Autorité}
\textbf{Faits.} L'Autorité de protection des données demande à NOUN PAY~:
la liste des données utilisées par son modèle de notation de crédit, le
lieu où elles sont hébergées, la base juridique du traitement, et la
manière dont un client refusé peut contester la décision.

\medskip
\textbf{Constatation.} NOUN PAY ne dispose d'aucun inventaire de ses
traitements. Ses données sont hébergées dans un centre de données situé
hors d'Afrique, en vertu d'un contrat type du fournisseur. Aucune
procédure ne permet à un client de contester un refus.

\medskip
\textbf{Analyse} (confiance moyenne, sous réserve de la lecture du texte
de la loi). Trois manquements probables~: l'absence d'inventaire, un
transfert hors du territoire non examiné au regard de la loi, et
l'absence de recours contre une décision automatisée. Le
Chapitre~\ref{chap:donnees-ia} reprendra le cas sous l'angle de la
gouvernance de l'intelligence artificielle.
\end{boxscenario}

% ============================================================================
\section{La Réglementation Bancaire de la CEMAC}
\label{sec:c10-cobac}
% ============================================================================

\subsection{Le règlement COBAC R-2016/04}

Les établissements de crédit de la zone CEMAC sont soumis à la
réglementation de la Commission bancaire de l'Afrique centrale. Le
règlement COBAC R-2016/04, relatif au contrôle interne dans les
établissements de crédit et les holdings financières, est le texte de
référence pour la gouvernance des SI d'une banque~\cite{cobac2016r04}.

\begin{boxjuris}
\textbf{Règlement COBAC R-2016/04, articles~1 et~2} (texte relevé sur une
base réglementaire en ligne)~:

\medskip
Article~1~: \textit{« Le présent règlement fixe les exigences minimales
relatives au système de contrôle interne que les établissements assujettis
doivent mettre en œuvre. »}

\medskip
Article~2, définition du risque opérationnel~: \textit{« risque de pertes
résultant de l'inadéquation ou de la défaillance des processus internes,
des personnes et des systèmes, ou d'évènements externes. »} Le même
article précise que ce risque inclut le risque juridique, et exclut les
risques stratégiques et de réputation.

\medskip
\textbf{Analyse.} La défaillance d'un système informatique est, par
définition, un risque opérationnel au sens du règlement. La panne de la
BANQUE SANAGA relève donc du dispositif de contrôle interne exigé par la
COBAC, et non d'un simple incident technique. Selon les présentations du
règlement, celui-ci impose en outre un dispositif d'identification, de
mesure et de surveillance de chaque risque, appuyé sur une cartographie
des risques, et compte un système de traitement de l'information parmi
les composantes du dispositif de contrôle interne.
\end{boxjuris}
% A COLLATIONNER avec la version officielle publiée par le SG-COBAC
% (sgcobac.org) : texte exact des articles 1 et 2, contenu des articles 3,
% 10 et 12.

\subsection{Ce que cela implique}

\begin{boxmetier}{Responsable GRC}
\textbf{Rédiger la partie SI du rapport à la COBAC}
\begin{enumerate}
    \item Qualifier la panne~: incident de risque opérationnel, avec sa
          chronologie et son coût (Chapitre~\ref{chap:prologue}).
    \item Situer la défaillance dans le dispositif de contrôle interne~:
          quels contrôles existaient, lesquels n'ont pas fonctionné.
    \item Relier le risque informatique à la cartographie des risques de la
          banque, avec l'échelle commune (Chapitre~\ref{chap:normes}).
    \item Présenter le plan d'action, objectif par objectif, à partir de la
          feuille de route COBIT (Chapitre~\ref{chap:cobit-capacite}).
\end{enumerate}
\end{boxmetier}

% ============================================================================
\section{Le Droit des Sociétés OHADA}
\label{sec:c10-ohada}
% ============================================================================

\subsection{La responsabilité des dirigeants}

L'Acte uniforme OHADA relatif au droit des sociétés commerciales et du
groupement d'intérêt économique, révisé le 30~janvier~2014~\cite{ohada2014auscgie},
organise la responsabilité civile des dirigeants sociaux. Selon la
doctrine, ses articles~161 et suivants posent le régime général~: le
dirigeant répond des fautes qu'il commet dans l'exercice de ses
fonctions, individuellement ou solidairement selon sa participation, et
peut être poursuivi par la société (action sociale) ou par les tiers
lésés. Pour la société anonyme, l'article~740 rend les administrateurs
responsables, individuellement ou solidairement, envers la société ou
envers les tiers.

\begin{boiteRemarque}{Un défaut de gouvernance des SI peut-il être une faute de gestion~?}
La question n'a pas, à notre connaissance, été tranchée par une
jurisprudence publiée en zone OHADA. Mais le raisonnement est
accessible~: si une banque soumise à des obligations précises (contrôle
interne, sécurité des systèmes) subit une perte majeure parce que ses
dirigeants ont maintenu une bascule malgré une alerte écrite, sans
décision formelle d'acceptation du risque, la qualification de faute de
gestion pourra être discutée. La trace écrite des décisions
(Chapitre~\ref{chap:prologue}) n'est donc pas une formalité~: elle est la
première défense d'un dirigeant diligent, et la première pièce d'un
dossier contre un dirigeant négligent.
\end{boiteRemarque}

\subsection{Le comité d'audit et le contrôle interne}

Le Chapitre~\ref{chap:gouv-entreprise} a présenté l'article~829-1 de l'Acte
uniforme révisé, qui rend le comité d'audit obligatoire dans les sociétés
cotées et le charge du suivi de l'efficacité des systèmes de contrôle
interne et de gestion des risques. Pour les entreprises et
établissements publics, les lois n°~2017/010 et 2017/011 et leur décret
d'application organisent le conseil d'administration et la tutelle~: la
gouvernance des SI de l'ANPP doit s'inscrire dans ce cadre, et dans celui
du régime financier de l'État (Chapitre~\ref{chap:alignement}).

% ============================================================================
\section{Le Cadre Africain}
\label{sec:c10-afrique}
% ============================================================================

\subsection{La Convention de Malabo}

La Convention de l'Union africaine sur la cybersécurité et la protection
des données à caractère personnel, adoptée à Malabo le 27~juin~2014, est
entrée en vigueur le 8~juin~2023, trente jours après le dépôt du
quinzième instrument de ratification. Elle couvre trois domaines~: les
transactions électroniques, la protection des données personnelles, et
la cybersécurité et la lutte contre la cybercriminalité. Elle oblige les
États parties à se doter d'une législation et d'autorités dans ces
domaines.

\begin{boxalert}
\textbf{À vérifier.} Nous n'avons pas trouvé de source établissant que le
Cameroun a ratifié la Convention de Malabo. Avant de l'invoquer comme
source de droit applicable au Cameroun, consultez la liste officielle des
signatures et ratifications tenue par l'Union africaine. Elle reste en
tout état de cause une référence de politique publique pour le
continent.
\end{boxalert}

\subsection{Une dynamique continentale}

Le Cameroun n'est pas seul. Plusieurs États africains se sont dotés de
lois de protection des données personnelles depuis la fin des années
2000, par exemple le Sénégal (loi n°~2008-12 du 25~janvier~2008) ou la
Côte d'Ivoire (loi n°~2013-450 du 19~juin~2013). Pour une organisation qui
opère dans plusieurs pays de la sous-région, comme une banque de la zone
CEMAC ou une entreprise de paiement mobile, la gouvernance des SI doit
tenir un registre des obligations pays par pays.

% ============================================================================
\section{Les Erreurs de Qualification}
\label{sec:c10-erreurs}
% ============================================================================

\begin{tableaucours}{L{4.6cm}L{7.6cm}}{%
    \tableauHeader{Erreur} &
    \tableauHeader{Pourquoi c'est une erreur}
}
« La panne est un incident technique, pas un sujet réglementaire. » & Pour
une banque, la défaillance d'un système est un risque opérationnel au sens
du règlement COBAC R-2016/04. \\
« Nous sommes certifiés ISO~27001, l'audit de l'ANTIC est inutile. » & La
certification est volontaire~; l'audit prévu par la loi n°~2010/012 est
obligatoire. \\
« Nos données sont chez un prestataire, c'est lui le responsable. » &
L'externalisation ne transfère pas la responsabilité du traitement ni
celle de la sécurité (loi n°~2010/012, art.~26~; loi n°~2024/017). \\
« Nos serveurs sont à l'étranger, la loi camerounaise ne s'applique
pas. » & Le transfert hors du territoire est précisément encadré par la
loi n°~2024/017. \\
« Seule la société est poursuivie, pas ses dirigeants. » & La
responsabilité de la personne morale n'exclut pas celle des personnes
physiques (loi n°~2010/012, art.~64~; AUSCGIE). \\
« La Convention de Malabo s'applique directement. » & Son applicabilité
au Cameroun suppose une ratification qui reste à établir. \\
\end{tableaucours}
\vspace{-0.8em}
\begin{center}{\small\itshape Tableau~10.4. Erreurs de qualification
fréquentes}\end{center}

\begin{boxprocedure}
\textbf{Le registre des obligations légales.} Pour démontrer la
conformité (objectif MEA03), l'organisation tient un registre qui, pour
chaque obligation, indique~:
\begin{itemize}
    \item le texte et l'article, avec la date de la version consultée~;
    \item l'obligation, reformulée en termes opérationnels~;
    \item le processus et le propriétaire qui la mettent en œuvre~;
    \item la preuve de conformité (document, enregistrement, rapport
          d'audit)~;
    \item la date de la dernière vérification et son résultat.
\end{itemize}
Le registre est revu au moins une fois par an, et à chaque nouveau texte.
\end{boxprocedure}

\begin{boxmetier}{Auditeur SI}
\textbf{Auditer la conformité légale d'un SI}
\begin{enumerate}
    \item Demander le registre des obligations, et vérifier qu'il couvre au
          moins les textes du Tableau~10.1 applicables à l'organisation.
    \item Pour un échantillon d'obligations, demander la preuve de
          conformité.
    \item Vérifier la conservation des journaux (durée, intégrité) au
          regard de l'article~29 de la loi n°~2010/012.
    \item Vérifier que les contrats des sous-traitants traitent de la
          sécurité, de la localisation des données et de l'audit.
\end{enumerate}
\end{boxmetier}

% ============================================================================
\separateur
% ============================================================================

\begin{boxresume}
\textbf{Ce qu'il faut retenir de ce chapitre~:}
\begin{itemize}
    \item \textbf{La loi n°~2010/012} confie la sécurité électronique à
          l'ANTIC et impose aux exploitants de systèmes d'information des
          mesures de sécurité et des dispositifs normalisés de gestion des
          risques (art.~26), la conservation des données de connexion
          pendant dix ans (art.~29), la révision des dispositifs (art.~30)
          et des audits obligatoires (art.~13 et~32).
    \item \textbf{La responsabilité est pénale aussi}~: la personne morale
          répond des infractions commises par ses organes, sans exonérer
          les personnes physiques (art.~64).
    \item \textbf{La loi n°~2024/017}, applicable depuis juin~2026, organise
          la protection des données personnelles et crée une autorité
          dédiée~: inventaire, sécurité, sous-traitants, localisation,
          preuve.
    \item \textbf{Pour les banques}, la défaillance d'un système est un
          risque opérationnel au sens du règlement COBAC R-2016/04.
    \item \textbf{Le droit OHADA} organise la responsabilité civile des
          dirigeants et le comité d'audit des sociétés cotées~: la trace
          écrite des décisions est la première défense.
    \item \textbf{La Convention de Malabo} est en vigueur depuis 2023~;
          son applicabilité au Cameroun est à vérifier.
    \item \textbf{Ne citer que ce qu'on a vérifié}~: un rapport au
          régulateur se fonde sur le texte publié, pas sur un résumé.
\end{itemize}
\end{boxresume}

\bigskip

\begin{boxexo}[title={\ding{45}\quad Exercice~10.1. Qui est compétent~? \niveaubase}]
Pour chacune des situations suivantes, indiquez l'autorité principalement
compétente (Tableau~10.3) et le texte en cause.
\begin{enumerate}[(a)]
    \item Une banque n'a pas fait réaliser l'audit de sécurité de ses
          systèmes depuis trois ans.
    \item Un client de NOUN PAY conteste l'usage de ses données par le
          modèle de notation.
    \item Le dispositif de contrôle interne d'une banque n'a pas détecté
          une migration non autorisée.
    \item Un établissement public a acquis un logiciel hors de tout
          programme budgétaire.
\end{enumerate}
\end{boxexo}

\begin{boxexo}[title={\ding{45}\quad Exercice~10.2. Le registre de KRIBILOG \niveaumoyen}]
KRIBILOG SA exploite un système de suivi des conteneurs qui enregistre les
données des chauffeurs (nom, numéro de permis, téléphone) et les journaux
d'accès des clients au portail web. L'exploitation est confiée à NETPORT
SERVICES, qui héberge une partie des données à l'étranger.
\begin{enumerate}[(a)]
    \item Établissez les cinq premières lignes du registre des obligations
          légales de KRIBILOG, sur le modèle de l'exigence de contrôle de
          la section~\ref{sec:c10-erreurs}.
    \item Pour chacune, indiquez si le texte a pu être vérifié, et sinon,
          où le vérifier.
    \item Quelles clauses faut-il ajouter au contrat de NETPORT SERVICES~?
\end{enumerate}
\end{boxexo}

\begin{boxexo}[title={\ding{45}\quad Exercice~10.3. La note au conseil de la BANQUE SANAGA \niveauexpert}]
Rédigez, à l'attention du conseil d'administration, la note juridique
évoquée en ouverture du chapitre.
\begin{enumerate}[(a)]
    \item Qualifiez la panne au regard du règlement COBAC R-2016/04 et de la
          loi n°~2010/012, en citant uniquement les articles vérifiés.
    \item Identifiez les obligations probablement méconnues, en séparant
          constatations et analyses, et en donnant un niveau de confiance.
    \item Discutez l'exposition éventuelle des dirigeants au regard du
          droit OHADA, en signalant ce qui reste incertain.
    \item Proposez trois mesures immédiates pour réduire le risque
          juridique de la banque.
\end{enumerate}
\end{boxexo}

\bigskip

\begin{boxinfo}
\textbf{Transition.} Les obligations ne valent que si quelqu'un vérifie
qu'elles sont tenues. Le Chapitre~\ref{chap:audit} présente l'audit, le
contrôle et l'assurance de la gouvernance des SI~: le modèle des trois
lignes, les normes de l'audit interne, l'audit public selon l'INTOSAI, et
le test des contrôles généraux informatiques.
\end{boxinfo}
